% ===== 第四部分：Fourier 分析 / 第2章：Fourier级数的收敛性 =====
% !TEX root = ../../main.tex
\setchapterimage{figure/P4/P4C2封面图.png}
\setchapterintro{本章集中讨论 Fourier 级数在不同意义下的收敛问题。我们先研究均方收敛与逐点收敛，说明正交展开如何逼近原函数；随后构造连续函数的 Fourier 级数仍可能发散的反例，并在最后整理函数正则性与 Fourier 系数衰减速度之间的关系。}
\chapter{Fourier级数的收敛性}
% 第1节：Fourier级数的均方收敛
% !TEX root = ../../../main.tex
% 所属：第四部分：Fourier 分析 / 第2章：Fourier级数的收敛性
% 节：2.1 Fourier级数的均方收敛
% 正文按原文件顺序保存；由 ../chapter2.tex 调用。
\section{Fourier级数的均方收敛}

\subsection*{预备知识}

默认读者已经对内积与范数,酉空间有初步了解.下面先介绍本节中至关重要的无穷维空间 $\ell^2(\Z)$.


\begin{definition}
	$\ell^2(\Z)$ 是所有复的双向无穷序列
	\[
	A=(\ldots,a_{-n},\ldots,a_{-1},a_0,a_1,\ldots,a_n,\ldots)
	\]
	所成的集合,其中
	\[
	\sum_{n\in\Z}|a_n|^2<\infty.
	\]
	向量的加法与数乘均按分量定义.若
	\[
	A=(a_n)_{n\in\Z},\qquad B=(b_n)_{n\in\Z},
	\]
	则定义内积
	\[
	(A,B)=\sum_{n\in\Z}a_n\overline{b_n}.
	\]
	该级数绝对收敛,并由此定义范数
	\[
	\norm{A}=(A,A)^{1/2}
	=\left(\sum_{n\in\Z}|a_n|^2\right)^{1/2}.
	\]
	
\end{definition}

下面证明 $\ell^2(\Z)$ 中的三角不等式与 Cauchy--Schwarz 不等式.

\begin{proposition}
	对任意 $A,B\in\ell^2(\Z)$,有
	\begin{enumerate}
		\item $\norm{A+B}\leq \norm{A}+\norm{B}$;
		\item $|(A,B)|\leq \norm{A}\norm{B}$.
	\end{enumerate}
\end{proposition}

\begin{proof}
	令
	\[
	A_N=(\ldots,0,a_{-N},\ldots,a_N,0,\ldots),
	\qquad
	B_N=(\ldots,0,b_{-N},\ldots,b_N,0,\ldots).
	\]
	在有限维空间中,三角不等式与 Cauchy--Schwarz 不等式成立,故
	\[
	\norm{A_N+B_N}\leq \norm{A_N}+\norm{B_N}
	\leq \norm{A}+\norm{B}.
	\]
	令 $N\to\infty$,得到
	\[
	\norm{A+B}\leq \norm{A}+\norm{B}.
	\]
	另一方面,
	\[
	|(A_N,B_N)|\leq \norm{A_N}\norm{B_N}
	\leq \norm{A}\norm{B}.
	\]
	同时,
	\[
	\sum_{|n|\leq N}|a_n b_n|
	\leq
	\left(\sum_{|n|\leq N}|a_n|^2\right)^{1/2}
	\left(\sum_{|n|\leq N}|b_n|^2\right)^{1/2}
	\leq \norm{A}\norm{B},
	\]
	所以 $\sum_{n\in\Z}a_n\overline{b_n}$ 绝对收敛.再令 $N\to\infty$,即得
	\[
	|(A,B)|\leq \norm{A}\norm{B}.
	\]
\end{proof}
在 $\R^d$、$\C^d$ 与 $\ell^2(\Z)$ 中,内积和范数满足下面两个重要性质：
\begin{enumerate}
	\item 内积是正定的,即 $\norm{X}=0$ 当且仅当 $X=0$;
	\item 空间是完备的,即每个 Cauchy 列都收敛,且其极限仍属于该空间.
\end{enumerate}
\begin{definition}
	满足上述两个条件的内积空间称为 \textbf{Hilbert 空间}.若上面的两个条件中有任意的一个不满足,则空间成为\textbf{准Hilbert 空间}
\end{definition}
因此,\allowbreak $\R^d$ 与 $\C^d$ 是有限维 Hilbert 空间,而 $\ell^2(\Z)$ 是无穷维 Hilbert 空间.

事实上,若 $\{A^{(k)}\}_{k\geq1}$ 是 $\ell^2(\Z)$ 中的 Cauchy 列,其中
\[
A^{(k)}=(a_n^{(k)})_{n\in\Z},
\]
则每个坐标列 $\{a_n^{(k)}\}_{k\geq1}$ 都是 $\C$ 中的 Cauchy 列.记其极限为 $a_n$,并令 $A=(a_n)_{n\in\Z}$.对充分大的 $k,m$ 及任意 $N$,有
\[
\sum_{|n|\leq N}|a_n^{(k)}-a_n^{(m)}|^2<\eps^2.
\]
先令 $m\to\infty$,再令 $N\to\infty$,便得到 $A\in\ell^2(\Z)$ 且 $\norm{A^{(k)}-A}\leq\eps$.因此 $\ell^2(\Z)$ 完备.

现在给出本节中另一个重要的函数空间.可以证明这是一个准Hilbert 空间


\begin{definition}
	令 $\mathcal R$ 表示区间 $[0,2\pi]$ 上复值 Riemann 可积函数的集合,等价地,也可将其看作圆周上的 Riemann 可积函数空间.函数的加法与数乘按点定义,并定义
	\[
	(f,g)=\frac1{2\pi}\int_0^{2\pi}f(\theta)\overline{g(\theta)}\dd\theta.
	\]
	相应的范数为
	\[
	\norm{f}
	=\left(\frac1{2\pi}\int_0^{2\pi}|f(\theta)|^2\dd\theta\right)^{1/2}.
	\]
\end{definition}

首先证明这一函数空间中的 Cauchy--Schwarz 不等式.对任意实数 $A,B$ 有$2AB\leq A^2+B^2$.令 $A=\lambda^{1/2}|f(\theta)|$、$B=\lambda^{-1/2}|g(\theta)|$,其中 $\lambda>0$,则
\[
|f(\theta)g(\theta)|
\leq \frac12\left(\lambda|f(\theta)|^2+\lambda^{-1}|g(\theta)|^2\right).
\]
对 $\theta$ 积分可得
\[
|(f,g)|
\leq \frac12\left(\lambda\norm{f}^2+\lambda^{-1}\norm{g}^2\right).
\]
当 $\norm{f}\neq0$ 且 $\norm{g}\neq0$ 时,取 $\lambda=\norm{g}/\norm{f}$,便有
\[
|(f,g)|\leq \norm{f}\norm{g}.
\]
其余情形显然成立.三角不等式是该不等式的直接推论.

不过,在 $\mathcal R$ 中,正定性并不成立.例如,令 $h(0)=1$,而在其余点令 $h(\theta)=0$,则 $h\neq0$ 但 $\norm{h}=0$.更本质地,这个空间也不是完备的.为说明这一点,定义
\[
f(\theta)=
\begin{cases}
	0,&\theta=0,\\
	\log(1/\theta),&0<\theta\leq2\pi.
\end{cases}
\]
由于 $f$ 在 $0$ 附近无界,故 $f\notin\mathcal R$.进一步定义
\[
f_n(\theta)=
\begin{cases}
	0,&0\leq\theta\leq 1/n,\\
	f(\theta),&1/n<\theta\leq2\pi.
\end{cases}
\]
每个 $f_n$ 都属于 $\mathcal R$.又因为 $\log^2(1/\theta)$ 在 $0$ 附近可积,所以
\[
\norm{f_m-f_n}^2
=\frac1{2\pi}\int_{\min\{1/m,1/n\}}^{\max\{1/m,1/n\}}
\log^2(1/\theta)\dd\theta=\theta(\log\theta)^2-2\theta\log\theta+2\theta\bigg|_{\min\{1/m,1/n\}}^{\max\{1/m,1/n\}}\longrightarrow0.
\]
因此 $\{f_n\}$ 是 $\mathcal R$ 中的 Cauchy 列,但它不可能收敛到 $\mathcal R$ 中的函数.事实上,$f_n$ 在上述范数下收敛到 $f$,而任何与 $f$ 几乎处处相等的函数都不可能是有界的 Riemann 可积函数.

\subsection*{均方收敛与 Parseval 恒等式}

对每个整数 $n$,令$e_n(\theta)=e^{in\theta}.$则 $\{e_n\}_{n\in\Z}$ 是标准正交族,即
\[
(e_n,e_m)=
\begin{cases}
	1,&n=m,\\
	0,&n\neq m.
\end{cases}
\]
若 $f$ 是圆周上的 Riemann 可积函数,其 Fourier 系数记为
\[
a_n=\widehat f(n)=\frac1{2\pi}\int_0^{2\pi}f(\theta)e^{-in\theta}\dd\theta,
\]
则
\[
a_n=(f,e_n),
\qquad
S_N(f)=\sum_{|n|\leq N}a_ne_n.
\]
由标准正交性可知,对任意 $|n|\leq N$,
\[
\left(f-\sum_{|k|\leq N}a_ke_k,e_n\right)=0.
\]
因此,对任意复数 $b_n$,有
\begin{equation}\label{eq:orthogonal-decomposition}
	\left(f-\sum_{|n|\leq N}a_ne_n,\sum_{|n|\leq N}b_ne_n\right)=0.
\end{equation}
取 $b_n=a_n$,由勾股定理得到
\begin{equation}\label{eq:finite-parseval}
	\norm{f}^2
	=\norm{f-S_N(f)}^2+
	\sum_{|n|\leq N}|a_n|^2.
\end{equation}

\begin{lemma}
	\textbf{（最佳逼近）}
	若 $f$ 是圆周上的 Riemann 可积函数,其 Fourier 系数为 $a_n$,则对任意复数 $c_n$,
	\[
	\norm{f-S_N(f)}
	\leq
	\norm{f-\sum_{|n|\leq N}c_ne_n}.
	\]
	并且等号成立当且仅当对所有 $|n|\leq N$ 都有 $c_n=a_n$.
\end{lemma}

\begin{proof}
	由
	\[
	f-\sum_{|n|\leq N}c_ne_n
	=f-S_N(f)+\sum_{|n|\leq N}(a_n-c_n)e_n
	\]
	以及式 \eqref{eq:orthogonal-decomposition},可得
	\[
	\norm{f-\sum_{|n|\leq N}c_ne_n}^2
	=\norm{f-S_N(f)}^2+
	\sum_{|n|\leq N}|a_n-c_n|^2.
	\]
	结论立即成立.
\end{proof}

\begin{theorem}
	\textbf{（均方收敛与 Parseval 恒等式）}
	设 $f$ 是圆周上的 Riemann 可积函数,且
	\[
	f(\theta)\sim\sum_{n=-\infty}^{\infty}a_ne^{in\theta}.
	\]
	则：
	\begin{enumerate}
		\item 当 $N\to\infty$ 时,$f$ 的 Fourier 级数均方收敛于 $f$,即
		\[
		\norm{f-S_N(f)}\longrightarrow0;
		\]
		\item Parseval 恒等式成立：
		\[
		\sum_{n=-\infty}^{\infty}|a_n|^2
		=\frac1{2\pi}\int_0^{2\pi}|f(\theta)|^2\dd\theta
		=\norm{f}^2.
		\]
	\end{enumerate}
\end{theorem}

\begin{proof}
	先设 $f$ 在圆周上连续.给定 $\eps>0$,由 Stone--Weierstrass 定理,存在次数为 $M$ 的三角多项式 $P$,使得
	\[
	|f(\theta)-P(\theta)|<\eps,
	\qquad \theta\in[0,2\pi].
	\]
	于是 $\norm{f-P}<\eps$.当 $N\geq M$ 时,由最佳逼近引理,
	\[
	\norm{f-S_N(f)}\leq\norm{f-P}<\eps.
	\]
	因此连续函数的 Fourier 级数均方收敛于原函数.
	
	现在设 $f$ 仅为 Riemann 可积函数.令 $B\geq\sup_{\theta}|f(\theta)|$.由 Part1中的\ref{C在R中一致有界稠密},可取圆周上的连续函数 $g$,使得 $|g|\leq B$ 且 $\displaystyle\int_0^{2\pi}|f-g|\dd\theta$ 任意小.由于 $|f-g|\leq2B$,
	\[
	\norm{f-g}^2
	=\frac1{2\pi}\int_0^{2\pi}|f-g|^2\dd\theta
	\leq\frac{B}{\pi}\int_0^{2\pi}|f-g|\dd\theta.
	\]
	故可使 $\norm{f-g}<\eps/2$.再由连续情形,存在三角多项式 $P$,使得 $\norm{g-P}<\eps/2$.于是
	\[
	\norm{f-P}\leq\norm{f-g}+\norm{g-P}<\eps.
	\]
	设 $P$ 的次数为 $M$,则当 $N\geq M$ 时,最佳逼近引理给出
	\[
	\norm{f-S_N(f)}\leq\norm{f-P}<\eps.
	\]
	这证明了均方收敛.
	
	最后,由式 \eqref{eq:finite-parseval},
	\[
	\sum_{|n|\leq N}|a_n|^2
	=\norm{f}^2-\norm{f-S_N(f)}^2.
	\]
	令 $N\to\infty$,便得到 Parseval 恒等式.
\end{proof}
在实际计算与应用中,我们经常采取Fourier展开的三角形式\[f(\theta)\sim \frac{a_0}{2}+\sum_{n=1}^{\infty}(a_n\cos n\theta+b_n\sin \theta)\]
其Parseval公式为
\[\frac{1}{2\pi}\int_{-\pi}^{\pi}|f(\theta)|^2\dd\theta=\frac{a_0^2}{4}+\sum_{n=1}^{\infty}(a_n^2+b_n^2)\]

式 \eqref{eq:finite-parseval} 还给出了更一般的结论.若 $\{u_n\}_{n\in\Z}$ 是任意标准正交族,且 $a_n=(f,u_n)$,则对每个 $N$ 有
\[
\sum_{|n|\leq N}|a_n|^2\leq\norm{f}^2.
\]
令 $N\to\infty$,得到熟知的 Bessel 不等式
\[
\sum_{n=-\infty}^{\infty}|a_n|^2\leq\norm{f}^2.
\]
等号成立的条件是该标准正交族对 $f$ 是完备的,即
\[
\norm{f-\sum_{|n|\leq N}a_nu_n}\longrightarrow0.
\]

\begin{lemma}
	\textbf{（Riemann--Lebesgue 引理）}
	如果 $f$ 是圆周上的 Riemann 可积函数,那么当 $|n|\to\infty$ 时,
	\[
	\widehat f(n)\longrightarrow0.
	\]
	等价地,当 $N\to\infty$ 时,
	\[
	\int_0^{2\pi}f(\theta)\sin(N\theta)\dd\theta\longrightarrow0,
	\qquad
	\int_0^{2\pi}f(\theta)\cos(N\theta)\dd\theta\longrightarrow0.
	\]
\end{lemma}

\begin{proof}
	给定 $\eps>0$,由上面的均方收敛结论,可取三角多项式 $P$,使得 $\norm{f-P}<\eps$.若 $|n|$ 大于 $P$ 的次数,则 $\widehat P(n)=0$,从而由 Cauchy--Schwarz 不等式,
	\[
	|\widehat f(n)|
	=|(f-P,e_n)|
	\leq\norm{f-P}\norm{e_n}
	<\eps.
	\]
	由于 $\norm{e_n}=1$,结论成立.正弦与余弦形式由
	\[
	e^{in\theta}=\cos(n\theta)+i\sin(n\theta)
	\]
	立即得到.
\end{proof}

\begin{lemma}\label{Parseval 等式的双线性形式}
	设 $F$ 和 $G$ 是圆周上的 Riemann 可积函数,且
	\[
	F(\theta)\sim\sum_{n=-\infty}^{\infty}a_ne^{in\theta},
	\qquad
	G(\theta)\sim\sum_{n=-\infty}^{\infty}b_ne^{in\theta}.
	\]
	那么
	\[\frac1{2\pi}\int_0^{2\pi}F(\theta)\overline{G(\theta)}\dd\theta
	=\sum_{n=-\infty}^{\infty}a_n\overline{b_n}.
	\]
	并且右端级数绝对收敛.
\end{lemma}

\begin{proof}
	由 Parseval 恒等式与复内积空间中的极化恒等式
	\[
	(F,G)=\frac14\left[
	\norm{F+G}^2-\norm{F-G}^2
	+i\bigl(\norm{F+iG}^2-\norm{F-iG}^2\bigr)
	\right],
	\]
	分别对 $F+G$、$F-G$、$F+iG$ 与 $F-iG$ 使用 Parseval 恒等式,逐项化简便得到
	\[
	(F,G)=\sum_{n=-\infty}^{\infty}a_n\overline{b_n}.
	\]
	此外,由 $\ell^2(\Z)$ 中的 Cauchy--Schwarz 不等式,
	\[
	\sum_{n=-\infty}^{\infty}|a_n\overline{b_n}|
	\leq
	\left(\sum_{n=-\infty}^{\infty}|a_n|^2\right)^{1/2}
	\left(\sum_{n=-\infty}^{\infty}|b_n|^2\right)^{1/2}
	=\norm{F}\norm{G}<\infty.
	\]
\end{proof}
利用 Parseval 恒等式,可以证明如下的 Wirtinger--Poincar\'e 不等式,它建立了函数的范数与其导数范数之间的关系.
\begin{proposition}
	\begin{enumerate}
		\item[(a)] 设 $f$ 是 $T$-周期、连续且分段 $C^1$ 的函数,并满足$\displaystyle\int_0^T f(t)\dd t=0.$
		
		那么有Wirtinger不等式
\[ \int_0^T \abs{f(t)}^2\dd t
			\leq
			\frac{T^2}{4\pi^2}
			\int_0^T \abs{f'(t)}^2\dd t. \]

		并且等号成立当且仅当$f(t)=A\sin\dfrac{2\pi t}{T}
			+B\cos\dfrac{2\pi t}{T}.$
		\item[(b)] 在 (a) 的条件下,若 $g$ 是 $C^1$ 的 $T$-周期函数,那么\[ \abs{\int_0^T \overline{f(t)}g(t)\dd t}^2
			\leq
			\frac{T^2}{4\pi^2}
			\left(\int_0^T\abs{f(t)}^2\dd t\right)
			\left(\int_0^T\abs{g'(t)}^2\dd t\right). \]
		\item[(c)] 设 $[a,b]$ 是紧区间,$f\in C^1([a,b])$ 且
		$f(a)=f(b)=0$.那么有Poincar\'e 不等式 
		
		\begin{equation}\label{eq:poincare-interval}
			\int_a^b\abs{f(t)}^2\dd t
			\leq
			\frac{(b-a)^2}{\pi^2}
			\int_a^b\abs{f'(t)}^2\dd t. 
		\end{equation}
		
		等号成立当且仅当$f(t)=A\sin(\pi\dfrac{t-a}{b-a})$,并且常数 $(b-a)^2/\pi^2$ 不能改进.
	\end{enumerate}
\end{proposition}

\begin{proof}
	\textbf{(a)} 对周期为 $T$ 的函数,其 Fourier 系数为
	\[
		\widehat f(n)
		=
		\frac1T\int_0^T f(t)e^{-2\pi int/T}\dd t,
		\qquad n\in\Z.
	\]
	由 $f$ 的平均值为零,得到$\widehat f(0)=0.$Parseval 恒等式给出
\[ \int_0^T\abs{f(t)}^2\dd t
		=
		T\sum_{n\in\Z}\abs{\widehat f(n)}^2
		=
		T\sum_{n\neq0}\abs{\widehat f(n)}^2. \]


	下面计算 $f'$ 的 Fourier 系数.由于 $f$ 连续且分段 $C^1$,可在各个光滑区间上分部积分；内部边界项因 $f$ 连续而相消.再利用 $f(T)=f(0)$ 可得
	\[\widehat{f'}(n)
		=\frac{2\pi in}{T}\widehat f(n)\]
	(在\ref{f(n)=O(1/n2)}中证明了)	因此,再次使用 Parseval 恒等式,得到
\[\int_0^T\abs{f'(t)}^2\dd t
		=
		T\sum_{n\in\Z}\abs{\widehat{f'}(n)}^2=
		\frac{4\pi^2}{T}
		\sum_{n\neq0}n^2\abs{\widehat f(n)}^2.\]
	因为 $n^2\geq1$ 对所有 $n\neq0$ 成立,所以由上面的推导
	\[\int_0^T\abs{f(t)}^2\dd t
		=T\sum_{n\neq0}\abs{\widehat f(n)}^2\leq T\sum_{n\neq0}n^2\abs{\widehat f(n)}^2=\frac{T^2}{4\pi^2}
		\int_0^T\abs{f'(t)}^2\dd t.\]
	得证

	等号成立当且仅当
	\[
		(n^2-1)\abs{\widehat f(n)}^2=0,
		\qquad n\neq0.
	\]
	也就是说,除 $n=\pm1$ 外,其余 Fourier 系数均为零.故
	\[
		f(t)=c_1e^{2\pi it/T}+c_{-1}e^{-2\pi it/T},
	\]
	等价地,
	\[
		f(t)=A\sin\frac{2\pi t}{T}
		+B\cos\frac{2\pi t}{T}.
	\]
	若 $f$ 为实值函数,则 $A,B\in\R$；一般复值情形下可取
	$A,B\in\C$.
	
	\textbf{(b)} 令
	\[
		g_0=\frac1T\int_0^T g(t)\dd t,
		\qquad
		h(t)=g(t)-g_0.
	\]
	则
	\[
		\int_0^T h(t)\dd t=0,
		\qquad h'(t)=g'(t).
	\]
	又因为 $\displaystyle\int_0^T f(t)\dd t=0$,所以
	\[
		\int_0^T \overline{f(t)}g(t)\dd t
		=
		\int_0^T \overline{f(t)}h(t)\dd t.
	\]
	由 Cauchy--Schwarz 不等式及 (a) 应用于 $h$,有
	\[
	\begin{aligned}
		\abs{\int_0^T\overline{f(t)}g(t)\dd t}^2
		&\leq
		\left(\int_0^T\abs{f(t)}^2\dd t\right)
		\left(\int_0^T\abs{h(t)}^2\dd t\right)\leq
		\frac{T^2}{4\pi^2}
		\left(\int_0^T\abs{f(t)}^2\dd t\right)
		\left(\int_0^T\abs{h'(t)}^2\dd t\right)\\
		&=
		\frac{T^2}{4\pi^2}
		\left(\int_0^T\abs{f(t)}^2\dd t\right)
		\left(\int_0^T\abs{g'(t)}^2\dd t\right).
	\end{aligned}
	\]
	\textbf{(c)} 记 $L=b-a$.先在区间 $[a-L,a+L]$ 上定义
	\[
		F(t)=
		\begin{cases}
			-f(2a-t),&a-L\leq t\leq a,\\
			f(t),&a\leq t\leq a+L=b,
		\end{cases}
	\]
	再将 $F$ 延拓为周期 $2L$ 的函数.由于 $f(a)=f(b)=0$,函数
	$F$ 连续且分段 $C^1$.它关于点 $a$ 为奇函数,因此在一个周期上的平均值为零：
	\[
		\int_{a-L}^{a+L}F(t)\dd t=0.
	\]
	同时,
	\[
		\int_{a-L}^{a+L}\abs{F(t)}^2\dd t
		=2\int_a^b\abs{f(t)}^2\dd t,
	\]
	以及
	\[
		\int_{a-L}^{a+L}\abs{F'(t)}^2\dd t
		=2\int_a^b\abs{f'(t)}^2\dd t.
	\]
	对周期为 $T=2L$ 的函数 $F$ 使用 (a),得到
	\[2\int_a^b\abs{f(t)}^2\dd t
		\leq
		\frac{(2L)^2}{4\pi^2}
		\,2\int_a^b\abs{f'(t)}^2\dd t=
		\frac{L^2}{\pi^2}
		\,2\int_a^b\abs{f'(t)}^2\dd t.	\]
	整理即得
	\[
		\int_a^b\abs{f(t)}^2\dd t
		\leq
		\frac{(b-a)^2}{\pi^2}
		\int_a^b\abs{f'(t)}^2\dd t.
	\]

	由 (a) 的等号条件,等号成立时,$F$ 只能包含周期 $2L$ 的第一频率.又因为 $F$ 关于 $a$ 为奇函数,余弦项必须消失.因此
	\[
		f(t)=A\sin\left(\frac{\pi(t-a)}{b-a}\right).
	\]
	反过来,这类函数确实使等号成立.

	最后,取非零函数
	\[
		f(t)=\sin\left(\frac{\pi(t-a)}{b-a}\right).
	\]
	它在 \eqref{eq:poincare-interval} 中恰好取等号.因此若将常数
	$(b-a)^2/\pi^2$ 换成任意更小的常数,不等式便会对该函数失效,故该常数不能改进.
\end{proof}

%
% 第2节：Fourier级数的逐点收敛
% !TEX root = ../../../main.tex
% 所属：第四部分：Fourier 分析 / 第2章：Fourier级数的收敛性
% 节：2.2 Fourier级数的逐点收敛
% 正文按原文件顺序保存；由 ../chapter2.tex 调用。
\section{Fourier级数的逐点收敛}

\begin{theorem}
	令 $f$ 是圆周上的可积函数,且在 $\theta_0$ 点可微.那么当 $N\to\infty$ 时,
	\[
	S_N(f)(\theta_0)\longrightarrow f(\theta_0).
	\]
\end{theorem}

\begin{proof}
	定义
	\[
	F(t)=
	\begin{cases}
		\displaystyle\frac{f(\theta_0-t)-f(\theta_0)}{t},
		& t\neq0,\ \abs{t}<\pi,\\[3mm]
		-f'(\theta_0),&t=0.
	\end{cases}
	\]
	首先,因为 $f$ 在 $\theta_0$ 点可微,所以 $F$ 在 $0$ 附近有界.其次,对充分小的 $\delta>0$,函数 $F$ 在
	\[
	[-\pi,-\delta]\cup[\delta,\pi]
	\]
	上可积；再由分段可积函数的拼接性质可知,$F$ 在整个区间 $[-\pi,\pi]$ 上可积.
	
	由
	\[
	S_N(f)(\theta_0)=(f\ast D_N)(\theta_0),
	\qquad
	\frac1{2\pi}\int_{-\pi}^{\pi}D_N(t)\dd t=1,
	\]
	可得
	\[
	\begin{aligned}
		S_N(f)(\theta_0)-f(\theta_0)
		&=\frac1{2\pi}\int_{-\pi}^{\pi}f(\theta_0-t)D_N(t)\dd t-f(\theta_0)\\
		&=\frac1{2\pi}\int_{-\pi}^{\pi}
		\bigl[f(\theta_0-t)-f(\theta_0)\bigr]D_N(t)\dd t\\
		&=\frac1{2\pi}\int_{-\pi}^{\pi}F(t)tD_N(t)\dd t.
	\end{aligned}
	\]
	再回顾 Dirichlet 核的公式
	\[
	tD_N(t)
	=\frac{t}{\sin(t/2)}
	\sin\bigl((N+\tfrac12)t\bigr),
	\]
	其中,$t/\sin(t/2)$ 在区间 $[-\pi,\pi]$ 上连续.又因为
	\[
	\sin\bigl((N+\dfrac12)t\bigr)
	=\sin(Nt)\cos(t/2)+\cos(Nt)\sin(t/2),
	\]
	所以
	\[
	\begin{aligned}
		S_N(f)(\theta_0)-f(\theta_0)
		={}&\frac1{2\pi}\int_{-\pi}^{\pi}
		F(t)\frac{t\cos(t/2)}{\sin(t/2)}\sin(Nt)\dd t\\
		&+\frac1{2\pi}\int_{-\pi}^{\pi}F(t)t\cos(Nt)\dd t.
	\end{aligned}
	\]
	上式中两个系数函数均为 Riemann 可积函数.对这两个积分应用 Riemann--Lebesgue 引理,便得
	\[
	S_N(f)(\theta_0)-f(\theta_0)\longrightarrow0.
	\]
\end{proof}

可以看出,上述证明中可微条件只用于保证函数
\[
F(t)=\frac{f(\theta_0-t)-f(\theta_0)}{t}
\]
在原点附近有界.定义
\[
F(t)=
\begin{cases}
	\dfrac{f(\theta_0-t)-f(\theta_0)}{t},
	& t\neq 0,\ |t|<\pi,\\[6pt]
	0,  & t=0.
\end{cases}
\]
因此,如果仅假设 $f$ 在 $\theta_0$ 点满足
Lipschitz 条件,定理的结论仍然成立.

事实上,设存在 $M>0$ 及 $\theta_0$ 的一个邻域,使得对该邻域内
任意的 $\theta$ 都有
\[
\abs{f(\theta)-f(\theta_0)}
\leq M\abs{\theta-\theta_0}.
\]
令 $\theta=\theta_0-t$,则当 $t$ 充分小时,
\[
\abs{F(t)}
=
\frac{\abs{f(\theta_0-t)-f(\theta_0)}}{\abs{t}}
\leq M.
\]
故 $F$ 在原点附近有界,其余证明与上述定理完全相同.
此时也称 $f$ 在 $\theta_0$ 点满足一阶 H\"older 条件.

\begin{corollary}
	假设 $f$ 和 $g$ 是定义在圆周上的两个可积函数,且对于某个 $\theta_0$,存在包含 $\theta_0$ 的开区间 $I$,满足
	\[
	f(\theta)=g(\theta),
	\qquad \theta\in I.
	\]
	那么当 $N\to\infty$ 时,
	\[
	S_N(f)(\theta_0)-S_N(g)(\theta_0)\longrightarrow0.
	\]
\end{corollary}

\begin{proof}
	令 $h=f-g$.因为 $h$ 在 $I$ 上恒为 $0$,所以 $h$ 在 $\theta_0$ 处可微,且 $h(\theta_0)=0$.由上一个定理,
	\[
	S_N(h)(\theta_0)\longrightarrow h(\theta_0)=0.
	\]
	再由 Fourier 部分和的线性性,
	\[
	S_N(h)(\theta_0)
	=S_N(f)(\theta_0)-S_N(g)(\theta_0),
	\]
	结论成立.
\end{proof}
%
% 第3节：一个反例——连续函数的 Fourier 级数可能发散
% !TEX root = ../../../main.tex
% 所属：第四部分：Fourier 分析 / 第2章：Fourier级数的收敛性
% 节：2.3 一个反例——连续函数的 Fourier 级数可能发散
% 正文按原文件顺序保存；由 ../chapter2.tex 调用。
\section{一个反例——连续函数的 Fourier 级数可能发散}


上一节说明,若函数在一点可微,则其 Fourier 级数在该点收敛到函数值.下面构造一个连续周期函数,使其 Fourier 级数在指定点发散.这表明,上一节中的可微条件不能直接削弱为连续条件.

\begin{proposition}
	存在圆周上的连续函数 $F$,使得其 Fourier 部分和在原点无界.更准确地说,存在严格递增的正整数列 $\{M_k\}$,满足
	\[
	\abs{S_{M_k}(F)(0)}\longrightarrow\infty.
	\]
	因此,$F$ 的 Fourier 级数在 $\theta=0$ 处发散.通过平移,也可以使发散发生在任意预先指定的点 $\theta_0$.
\end{proposition}

构造的关键是 Fourier 部分和中的\textbf{对称性破缺}.通常
\[
S_N(f)(\theta)=\sum_{-N}^{N}\widehat f(n)e^{in\theta}
\]
同时截取正频率与负频率,二者之间可能产生显著抵消.我们将先构造一个一致有界的对称频率块,再通过适当移动频率,使某个 Fourier 部分和恰好从中间切开该频率块.被截出的单边部分将在原点具有对数级增长.



考虑 $h$ ,圆周上的锯齿函数：它关于 $\theta$ 为奇函数,并且当 $0<\theta<\pi$ 时,
\[
h(\theta)=i(\pi-\theta).
\]
由前面的计算可知,
\begin{equation}\label{eq:sawtooth-fourier}
	h(\theta)\sim\sum_{n\neq0}\frac{e^{in\theta}}{n}.
\end{equation}
式 \eqref{eq:sawtooth-fourier} 中正、负频率关于原点对称.若只保留负频率,得到形式级数
\[
\sum_{n=-\infty}^{-1}\frac{e^{in\theta}}{n}.
\]
这个单边级数不可能是某个 Riemann 可积函数的 Fourier 级数.事实上,若它是某个函数 $\widetilde h$ 的 Fourier 级数,则 $\widetilde h$ 有界,且在原点的 Abel 平均满足
\[
\abs{A_r(\widetilde h)(0)}
=
\sum_{n=1}^{\infty}\frac{r^n}{n}
\longrightarrow\infty,
\qquad r\to1^-.
\]
另一方面,由 Abel 平均的卷积表示及 Poisson 核的正性,
\[
\abs{A_r(\widetilde h)(0)}
=
\abs{(\widetilde h\ast P_r)(0)}
\leq
\sup_{\theta}\abs{\widetilde h(\theta)},
\]
得到矛盾.

对每个正整数 $N$,定义三角多项式

	\[ f_N(\theta)=\sum_{1\leq\abs{n}\leq N}\frac{e^{in\theta}}{n},
	\qquad
	\widetilde f_N(\theta)=\sum_{-N\leq n\leq-1}\frac{e^{in\theta}}{n}. \]

在原点,
\[
\abs{\widetilde f_N(0)}
=
\sum_{n=1}^{N}\frac1n
\geq
\int_1^N\frac{\dd x}{x}
=
\log N.
\]
因此存在绝对常数 $c>0$,使得
\begin{equation}\label{eq:one-sided-log-growth}
	\abs{\widetilde f_N(0)}\geq c\log N.
\end{equation}
与此相反,\textbf{完整的对称和 $f_N$ 关于 $N$ 与 $\theta$ 一致有界}.为证明这一点,先给出一个简单引理.

\begin{lemma}
	设复数列 $\{c_n\}_{n\geq1}$ 的 Abel 平均
	\[
	A_r=\sum_{n=1}^{\infty}r^nc_n,
	\qquad 0<r<1,
	\]
	在 $r\to1^-$ 时有界.若 $c_n=O(1/n)$,则其部分和
	\[
	T_N=\sum_{n=1}^{N}c_n
	\]
	也是有界的.
\end{lemma}

\begin{proof}
	取 $r=1-1/N$.选择常数 $M>0$,使得对所有 $n$ 及充分接近 $1$ 的 $r$ 都有
	\[
	n\abs{c_n}\leq M,
	\qquad
	\abs{A_r}\leq M.
	\]
	于是
	\[
	S_N-A_r
	=
	\sum_{n=1}^{N}(1-r^n)c_n
	-
	\sum_{n=N+1}^{\infty}r^nc_n.
	\]
	利用
	\[
	1-r^n=(1-r)(1+r+\cdots+r^{n-1})\leq n(1-r),
	\]
	得到
	\[
	\begin{aligned}
		\abs{S_N-A_r}
		&\leq
		\sum_{n=1}^{N}\abs{c_n}(1-r^n)
		+
		\sum_{n=N+1}^{\infty}r^n\abs{c_n}\\
		&\leq
		M\sum_{n=1}^{N}(1-r)
		+
		\frac{M}{N}\sum_{n=N+1}^{\infty}r^n\\
		&\leq
		MN(1-r)+\frac{M}{N(1-r)}
		=2M.
	\end{aligned}
	\]
	故 $\abs{S_N}\leq3M$,从而部分和有界.
\end{proof}

将引理应用于
\[c_n=\frac{e^{in\theta}}{n}+\frac{e^{i(-n)\theta}}{-n}\]
则
\[
f_N(\theta)=\sum_{n=1}^{N}c_n(\theta).
\]
显然 $n\abs{c_n(\theta)}\leq2$.另一方面,相应的 Abel 平均正是锯齿函数 $h$ 的 Abel 平均,
\[
\sum_{n=1}^{\infty}r^nc_n(\theta)
=A_r(h)(\theta)
=(h\ast P_r)(\theta),
\]
故
\[
\abs{A_r(h)(\theta)}\leq\sup_{\varphi}\abs{h(\varphi)}.
\]
上述估计与 $r,\theta$ 无关.因此存在绝对常数 $C>0$,使得
\begin{equation}\label{eq:fN-uniform-bound}
	\abs{f_N(\theta)}\leq C,
	\qquad N\geq1,\quad\theta\in\R.
\end{equation}



定义
\begin{equation}\label{eq:shifted-blocks}
	P_N(\theta)=e^{i2N\theta}f_N(\theta),
	\qquad
	\widetilde P_N(\theta)=e^{i2N\theta}\widetilde f_N(\theta).
\end{equation}
乘以 $e^{i2N\theta}$ 会把所有频率向右移动 $2N$ 个单位.因此,$P_N$ 的非零 Fourier 系数只出现在
\[
N\leq n\leq3N,
\qquad n\neq2N.
\]
其中区间 $[N,2N-1]$ 对应原来的负频率,区间 $[2N+1,3N]$ 对应原来的正频率.

\begin{lemma}\label{lem:frequency-block-cut}
	对任意正整数 $M,N$,有
	\[
	S_M(P_N)=
	\begin{cases}
		P_N,&M\geq3N,\\
		\widetilde P_N,&M=2N,\\
		0,&M<N.
	\end{cases}
	\]
\end{lemma}

\begin{proof}
	当 $M<N$ 时,截断区间 $[-M,M]$ 不包含 $P_N$ 的任何频率,故 $S_M(P_N)=0$.当 $M\geq3N$ 时,全部频率都被保留,故 $S_M(P_N)=P_N$.最后,当 $M=2N$ 时,仅保留区间 $[N,2N-1]$ 中的频率,这恰好是 $\widetilde f_N$ 的频率向右移动 $2N$ 后所得的 $\widetilde P_N$.
\end{proof}

注意,由式 \eqref{eq:fN-uniform-bound},
\begin{equation}\label{eq:PN-uniform-bound}
	\abs{P_N(\theta)}=\abs{f_N(\theta)}\leq C,
\end{equation}
但由式 \eqref{eq:one-sided-log-growth},
\begin{equation}\label{eq:Ptilde-log-growth}
	\abs{\widetilde P_N(0)}
	=
	\abs{\widetilde f_N(0)}
	\geq c\log N.
\end{equation}
因此,$P_N$ 本身一致有界,而部分和算子 $S_{2N}$ 从频率块中截出的部分却在原点具有对数级增长.



\begin{proof}[命题的证明]
	选择正数列与正整数列
	\[
	\alpha_k=\frac1{k^2},
	\qquad
	N_k=3^{2^k}.
	\]
	则
	\[
	\sum_{k=1}^{\infty}\alpha_k<\infty,
	\qquad
	N_{k+1}>3N_k,
	\qquad
	\alpha_k\log N_k
	=
	\frac{2^k\log3}{k^2}
	\longrightarrow\infty.
	\]
	定义
	\begin{equation}\label{eq:continuous-divergent-example}
		F(\theta)=\sum_{k=1}^{\infty}\alpha_kP_{N_k}(\theta).
	\end{equation}
	由式 \eqref{eq:PN-uniform-bound} 与 Weierstrass 判别法,级数 \eqref{eq:continuous-divergent-example} 一致收敛.因此 $F$ 是圆周上的连续函数.
	
	由于 $N_{k+1}>3N_k$,不同 $P_{N_k}$ 的频率支集彼此分离.固定 $m$,考察阶数为 $2N_m$ 的 Fourier 部分和.由引理 \ref{lem:frequency-block-cut},当 $k>m$ 时,
	\[
	S_{2N_m}(P_{N_k})=0;
	\]
	当 $k<m$ 时,
	\[
	S_{2N_m}(P_{N_k})=P_{N_k};
	\]
	而当 $k=m$ 时,
	\[
	S_{2N_m}(P_{N_m})=\widetilde P_{N_m}.
	\]
	又因为
	\[
	P_{N_k}(0)=f_{N_k}(0)
	=
	\sum_{1\leq\abs{n}\leq N_k}\frac1n
	=0,
	\]
	所以
	\[
	S_{2N_m}(F)(0)
	=
	\alpha_m\widetilde P_{N_m}(0).
	\]
	由式 \eqref{eq:Ptilde-log-growth},
	\[
	\abs{S_{2N_m}(F)(0)}
	\geq
	c\alpha_m\log N_m
	=
	c\frac{2^m\log3}{m^2}
	\longrightarrow\infty.
	\]
	故 $F$ 的 Fourier 部分和在原点无界,从而其 Fourier 级数在原点发散.
	
	若希望发散发生在预先指定的点 $\theta_0$,只需令
	\[
	F_{\theta_0}(\theta)=F(\theta-\theta_0).
	\]
	此时
	\[
	S_N(F_{\theta_0})(\theta_0)=S_N(F)(0),
	\]
	结论成立.
\end{proof}
%
% 第4节：Fourier系数的衰减估计
% !TEX root = ../../../main.tex
% 所属：第四部分：Fourier 分析 / 第2章：Fourier级数的收敛性
% 节：2.4 Fourier系数的衰减估计
% 正文按原文件顺序保存；由 ../chapter2.tex 调用。
\section{Fourier系数的衰减估计}

\begin{proposition}

设 $f$ 是 $2\pi$-周期的 $C^k$ 函数,则
\[
\widehat f(n)=o\left(\frac1{\abs{n}^k}\right),
\qquad \abs{n}\to\infty.
\]
等价地,
\[
\abs{n}^k\widehat f(n)\longrightarrow0.
\]
\end{proposition}

\begin{proof}
对 $n\neq0$ 分部积分.由于 $f$ 是周期函数,
\[
f(\pi)=f(-\pi),
\qquad
 e^{-in\pi}=e^{in\pi}=(-1)^n,
\]
故边界项消失,从而
\[
\widehat f(n)
=\frac1{in}\widehat{f'}(n).
\]
周期函数的各阶导数仍是周期函数,因此重复分部积分 $k$ 次可得

\[ \widehat f(n)=\frac1{(in)^k}\widehat{f^{(k)}}(n). \]

于是
\[
\abs{n}^k\abs{\widehat f(n)}
=\abs{\widehat{f^{(k)}}(n)}.
\]
函数 $f^{(k)}$ 连续,因而 Riemann 可积.由 Riemann--Lebesgue 引理,
\[
\abs{n}^k\widehat f(n)=\widehat{f^{(k)}}(n)\longrightarrow0,
\qquad \abs{n}\to\infty.
\]

\end{proof}

\ref{f(n)=O(1/n^k)}只能直接给出
\[
\widehat f(n)=O\left(\frac1{\abs{n}^k}\right),
\]
而 Riemann--Lebesgue 引理进一步把 $O$ 改进为 $o$.

\begin{lemma}

设 $f$ 是 $2\pi$-周期函数,并且在 $[-\pi,\pi]$ 上 Riemann 可积.对每个 $n\neq0$,有

\[ \widehat f(n)
=-\frac1{2\pi}\int_{-\pi}^{\pi}
 f\left(x+\frac{\pi}{n}\right)e^{-inx}\dd x, \]

从而

\[ \widehat f(n)
=\frac1{4\pi}\int_{-\pi}^{\pi}
\left[f(x)-f\left(x+\frac{\pi}{n}\right)\right]
 e^{-inx}\dd x. \]

\end{lemma}

\begin{proof}
在积分
\[
\frac1{2\pi}\int_{-\pi}^{\pi}
 f\left(x+\frac{\pi}{n}\right)e^{-inx}\dd x
\]
中作变量代换 $y=x+\pi/n$,则
\[
e^{-inx}=e^{-in(y-\pi/n)}=-e^{-iny}.
\]
又因为 $f(y)e^{-iny}$ 是 $2\pi$-周期函数,所以它在任意长度为 $2\pi$ 的区间上的积分相同.因此
\[
\frac1{2\pi}\int_{-\pi}^{\pi}
 f\left(x+\frac{\pi}{n}\right)e^{-inx}\dd x
=-\widehat f(n),
\]
即得.再将该式与 Fourier 系数的定义相加,便得证
\end{proof}

\begin{proposition}
若 $f$ 满足阶数为 $\alpha$ 的 Hölder 条件
\[
\abs{f(x+h)-f(x)}\leq C\abs{h}^{\alpha},
\qquad 0<\alpha\leq1,
\]
则
\[
\widehat f(n)=O\left(\frac1{\abs{n}^{\alpha}}\right).
\]
更准确地说,对 $n\neq0$,
\[
\abs{\widehat f(n)}
\leq \frac{C\pi^{\alpha}}{2\abs{n}^{\alpha}}.
\]
\end{proposition}

\begin{proof}
由上面引理的第二个公式 以及 $\abs{e^{-inx}}=1$,
\[
\begin{aligned}
\abs{\widehat f(n)}
&=\frac1{4\pi}\int_{-\pi}^{\pi}
\abs{f(x)-f\left(x+\frac{\pi}{n}\right)}\dd x\\
&\leq\frac1{4\pi}\int_{-\pi}^{\pi}
C\left(\frac{\pi}{\abs{n}}\right)^{\alpha}\dd x
=\frac{C\pi^{\alpha}}{2\abs{n}^{\alpha}}.
\end{aligned}
\]
结论成立.
\end{proof}
上面的命题在 $\alpha=1$ 时只给出 $\widehat f(n)=O(1/\abs{n})$.由于调和级数发散,这一逐项估计尚不足以推出 Fourier 级数绝对收敛.下面利用 Parseval 恒等式对频率作二进分块,得到更强的结论.

\begin{proposition}
设 $f$ 是 $2\pi$-周期函数,并满足 Lipschitz 条件
\[
\abs{f(x)-f(y)}\leq K\abs{x-y}.
\]
则
\[
\sum_{n\in\Z}\abs{\widehat f(n)}<\infty.
\]
因此 $f$ 的 Fourier 级数绝对且一致收敛.

更一般地,若 $f$ 满足阶数 $\alpha>1/2$ 的 Hölder 条件,则其 Fourier 级数绝对收敛.这一结论称为 Bernstein 定理.
\end{proposition}

\begin{proof}
对 $h>0$ 定义
\[
g_h(x)=f(x+h)-f(x-h).
\]
由 Fourier 系数的平移性质,
\[
\widehat{g_h}(n)
=\left(e^{inh}-e^{-inh}\right)\widehat f(n)
=2i\sin(nh)\widehat f(n).
\]
Parseval 恒等式给出

\begin{equation}\label{eq:gh-parseval}
	\frac1{2\pi}\int_0^{2\pi}\abs{g_h(x)}^2\dd x
=4\sum_{n\in\Z}\abs{\sin(nh)}^2\abs{\widehat f(n)}^2.
\end{equation} 


先考虑 Lipschitz 情形.此时
\[
\abs{g_h(x)}
=\abs{f(x+h)-f(x-h)}
\leq2Kh.
\]
那么,

\[ \sum_{n\in\Z}\abs{\sin(nh)}^2\abs{\widehat f(n)}^2
\leq K^2h^2. \]


固定正整数 $p$,取$h=\dfrac{\pi}{2^{p+1}}.$当 $2^{p-1}<\abs{n}\leq2^p$ 时,
\[
\frac{\pi}{4}<\abs{n}h\leq\frac{\pi}{2},
\qquad
\abs{\sin(nh)}^2\geq\frac12.
\]
故\[ \sum_{2^{p-1}<\abs{n}\leq2^p}\abs{\widehat f(n)}^2
\leq\frac{K^2\pi^2}{2^{2p+1}}. \]

该频率块中共有 $2^p$ 个整数.由 Cauchy--Schwarz 不等式,
\[\sum_{2^{p-1}<\abs{n}\leq2^p}\abs{\widehat f(n)}\leq 2^{p/2}
\left(
\sum_{2^{p-1}<\abs{n}\leq2^p}\abs{\widehat f(n)}^2
\right)^{1/2}\leq \frac{K\pi}{2^{(p+1)/2}}.\]
对 $p$ 求和便得
\[
\sum_{n\in\Z}\abs{\widehat f(n)}<\infty.
\]
又因为$\abs{\widehat f(n)e^{inx}}=\abs{\widehat f(n)},$由 Weierstrass 判别法,Fourier 级数一致收敛.

现在设 $f$ 满足阶数为 $\alpha>1/2$ 的 Hölder 条件.则
\[
\abs{g_h(x)}\leq K(2h)^{\alpha}.
\]
利用式 \eqref{eq:gh-parseval},得到
\[
\sum_{n\in\Z}\abs{\sin(nh)}^2\abs{\widehat f(n)}^2
\leq K^2 2^{2\alpha-2}h^{2\alpha}.
\]
仍取 $h=\pi/2^{p+1}$,可得
\[
\sum_{2^{p-1}<\abs{n}\leq2^p}\abs{\widehat f(n)}^2
\leq K^2\pi^{2\alpha}2^{-2\alpha p-1}.
\]
再次应用 Cauchy--Schwarz 不等式,
\[
\sum_{2^{p-1}<\abs{n}\leq2^p}\abs{\widehat f(n)}
\leq\frac{K\pi^{\alpha}}{\sqrt2}
2^{-p(\alpha-1/2)}.
\]
由于 $\alpha>1/2$,右侧关于 $p$ 可求和,故
\[
\sum_{n\in\Z}\abs{\widehat f(n)}<\infty.
\]
Bernstein 定理得证.
\end{proof}
\begin{proposition}
若 $f$ 是 $[-\pi,\pi]$ 上的有界单调函数,则
\[
\widehat f(n)=O\left(\frac1{\abs{n}}\right).
\]
\end{proposition}

\begin{proof}
不妨设 $f$ 单调递增,且 $\abs{f(x)}\leq M$.首先,对任意区间 $[a,b]\subset[-\pi,\pi]$ 及 $n\neq0$,
\[
\widehat{\chi_{[a,b]}}(n)
=\frac1{2\pi}\int_a^b e^{-inx}\dd x
=\frac{e^{-ina}-e^{-inb}}{2\pi in},
\]
故
\begin{equation}\label{eq:interval-indicator-decay}
\abs{\widehat{\chi_{[a,b]}}(n)}
\leq\frac1{\pi\abs{n}}.
\end{equation}

设 $g$ 是单调递增的阶梯函数,
\[
g(x)=\sum_{k=1}^{N}\alpha_k\chi_{[a_k,a_{k+1}]}(x),
\]
其中
\[
-\pi=a_1<a_2<\cdots<a_{N+1}=\pi,
\qquad
-M\leq\alpha_1\leq\cdots\leq\alpha_N\leq M.
\]
忽略有限个端点处的取值,可将其写为
\[
g(x)=\alpha_1+
\sum_{k=1}^{N-1}(\alpha_{k+1}-\alpha_k)
\chi_{[a_{k+1},\pi]}(x).
\]
当 $n\neq0$ 时,常数项的 Fourier 系数为零.因此由式 \eqref{eq:interval-indicator-decay},
\[\abs{\widehat g(n)}
\leq\frac1{\pi\abs{n}}
\sum_{k=1}^{N-1}(\alpha_{k+1}-\alpha_k)=\frac{\alpha_N-\alpha_1}{\pi\abs{n}}
\leq\frac{2M}{\pi\abs{n}}.\]
这一估计与阶梯数 $N$ 无关.

最后用上述阶梯函数一致逼近 $f$.对正整数 $m$,定义
\[
g_m(x)=f(-\pi)+\frac1m
\left\lfloor m\bigl(f(x)-f(-\pi)\bigr)\right\rfloor.
\]
因为 $f$ 单调且有界,$g_m$ 是只取有限多个值的单调阶梯函数,并且
\[
0\leq f(x)-g_m(x)<\frac1m.
\]
同时 $g_m$ 的值介于 $f(-\pi)$ 与 $f(\pi)$ 之间,故前面的估计给出
\[
\abs{\widehat{g_m}(n)}\leq\frac{2M}{\pi\abs{n}}.
\]
又因为 $g_m\to f$ 一致收敛,对每个固定的 $n$,
\[
\abs{\widehat f(n)-\widehat{g_m}(n)}
\leq\norm{f-g_m}_{\infty}\longrightarrow0.
\]
令 $m\to\infty$,便得到
\[
\abs{\widehat f(n)}\leq\frac{2M}{\pi\abs{n}}.
\]
若 $f$ 单调递减,则对 $-f$ 应用上述结论即可.
\end{proof}



前面的结果表明,额外的限制条件会带来明确的衰减速度.然而,连续性本身并不能保证任何预先指定的统一速度.也就是说,连续函数的 Fourier 系数可以衰减得任意慢.

\begin{proposition}
设 $\{\eps_n\}_{n\geq1}$ 是任意非负实数列,且 $\eps_n\to0$.则存在连续的 $2\pi$-周期函数 $f$,使得
\[
\abs{\widehat f(n)}\geq\eps_n
\]
对无穷多个正整数 $n$ 成立.
\end{proposition}

\begin{proof}
因为 $\eps_n\to0$,可以递归选择严格递增的正整数列$n_1<n_2<\cdots$使得$\eps_{n_k}\leq2^{-k}.$于是
\[
\sum_{k=1}^{\infty}\eps_{n_k}<\infty.
\]
定义
\begin{equation}\label{eq:arbitrarily-slow-decay-function}
f(x)=\sum_{k=1}^{\infty}\eps_{n_k}e^{in_kx}.
\end{equation}
由 Weierstrass 判别法,式 \eqref{eq:arbitrarily-slow-decay-function} 在 $\R$ 上绝对且一致收敛.因此 $f$ 是连续的 $2\pi$-周期函数.

一致收敛允许逐项积分.由复指数函数的正交性,
\[
\begin{aligned}
\widehat f(n)
&=\sum_{k=1}^{\infty}\eps_{n_k}
\frac1{2\pi}\int_{-\pi}^{\pi}e^{i(n_k-n)x}\dd x\\
&=
\begin{cases}
\eps_{n_k},&n=n_k\text{ 对某个 }k,\\
0,&n\notin\{n_1,n_2,\ldots\}.
\end{cases}
\end{aligned}
\]
特别地,
\[
\abs{\widehat f(n_k)}=\eps_{n_k}
\]
对每个 $k$ 都成立,所以所需不等式对无穷多个 $n$ 成立.
\end{proof}

\textbf{注：}
若希望得到实值函数,可改取
\[
f(x)=2\sum_{k=1}^{\infty}\eps_{n_k}\cos(n_kx).
\]
此时
\[
\widehat f(n_k)=\widehat f(-n_k)=\eps_{n_k}.
\]
%
